% Options for packages loaded elsewhere
\PassOptionsToPackage{unicode}{hyperref}
\PassOptionsToPackage{hyphens}{url}
\documentclass[
]{article}
\usepackage{xcolor}
\usepackage[margin=1in]{geometry}
\usepackage{amsmath,amssymb}
\setcounter{secnumdepth}{-\maxdimen} % remove section numbering
\usepackage{iftex}
\ifPDFTeX
  \usepackage[T1]{fontenc}
  \usepackage[utf8]{inputenc}
  \usepackage{textcomp} % provide euro and other symbols
\else % if luatex or xetex
  \usepackage{unicode-math} % this also loads fontspec
  \defaultfontfeatures{Scale=MatchLowercase}
  \defaultfontfeatures[\rmfamily]{Ligatures=TeX,Scale=1}
\fi
\usepackage{lmodern}
\ifPDFTeX\else
  % xetex/luatex font selection
\fi
% Use upquote if available, for straight quotes in verbatim environments
\IfFileExists{upquote.sty}{\usepackage{upquote}}{}
\IfFileExists{microtype.sty}{% use microtype if available
  \usepackage[]{microtype}
  \UseMicrotypeSet[protrusion]{basicmath} % disable protrusion for tt fonts
}{}
\makeatletter
\@ifundefined{KOMAClassName}{% if non-KOMA class
  \IfFileExists{parskip.sty}{%
    \usepackage{parskip}
  }{% else
    \setlength{\parindent}{0pt}
    \setlength{\parskip}{6pt plus 2pt minus 1pt}}
}{% if KOMA class
  \KOMAoptions{parskip=half}}
\makeatother
\usepackage{color}
\usepackage{fancyvrb}
\newcommand{\VerbBar}{|}
\newcommand{\VERB}{\Verb[commandchars=\\\{\}]}
\DefineVerbatimEnvironment{Highlighting}{Verbatim}{commandchars=\\\{\}}
% Add ',fontsize=\small' for more characters per line
\usepackage{framed}
\definecolor{shadecolor}{RGB}{248,248,248}
\newenvironment{Shaded}{\begin{snugshade}}{\end{snugshade}}
\newcommand{\AlertTok}[1]{\textcolor[rgb]{0.94,0.16,0.16}{#1}}
\newcommand{\AnnotationTok}[1]{\textcolor[rgb]{0.56,0.35,0.01}{\textbf{\textit{#1}}}}
\newcommand{\AttributeTok}[1]{\textcolor[rgb]{0.13,0.29,0.53}{#1}}
\newcommand{\BaseNTok}[1]{\textcolor[rgb]{0.00,0.00,0.81}{#1}}
\newcommand{\BuiltInTok}[1]{#1}
\newcommand{\CharTok}[1]{\textcolor[rgb]{0.31,0.60,0.02}{#1}}
\newcommand{\CommentTok}[1]{\textcolor[rgb]{0.56,0.35,0.01}{\textit{#1}}}
\newcommand{\CommentVarTok}[1]{\textcolor[rgb]{0.56,0.35,0.01}{\textbf{\textit{#1}}}}
\newcommand{\ConstantTok}[1]{\textcolor[rgb]{0.56,0.35,0.01}{#1}}
\newcommand{\ControlFlowTok}[1]{\textcolor[rgb]{0.13,0.29,0.53}{\textbf{#1}}}
\newcommand{\DataTypeTok}[1]{\textcolor[rgb]{0.13,0.29,0.53}{#1}}
\newcommand{\DecValTok}[1]{\textcolor[rgb]{0.00,0.00,0.81}{#1}}
\newcommand{\DocumentationTok}[1]{\textcolor[rgb]{0.56,0.35,0.01}{\textbf{\textit{#1}}}}
\newcommand{\ErrorTok}[1]{\textcolor[rgb]{0.64,0.00,0.00}{\textbf{#1}}}
\newcommand{\ExtensionTok}[1]{#1}
\newcommand{\FloatTok}[1]{\textcolor[rgb]{0.00,0.00,0.81}{#1}}
\newcommand{\FunctionTok}[1]{\textcolor[rgb]{0.13,0.29,0.53}{\textbf{#1}}}
\newcommand{\ImportTok}[1]{#1}
\newcommand{\InformationTok}[1]{\textcolor[rgb]{0.56,0.35,0.01}{\textbf{\textit{#1}}}}
\newcommand{\KeywordTok}[1]{\textcolor[rgb]{0.13,0.29,0.53}{\textbf{#1}}}
\newcommand{\NormalTok}[1]{#1}
\newcommand{\OperatorTok}[1]{\textcolor[rgb]{0.81,0.36,0.00}{\textbf{#1}}}
\newcommand{\OtherTok}[1]{\textcolor[rgb]{0.56,0.35,0.01}{#1}}
\newcommand{\PreprocessorTok}[1]{\textcolor[rgb]{0.56,0.35,0.01}{\textit{#1}}}
\newcommand{\RegionMarkerTok}[1]{#1}
\newcommand{\SpecialCharTok}[1]{\textcolor[rgb]{0.81,0.36,0.00}{\textbf{#1}}}
\newcommand{\SpecialStringTok}[1]{\textcolor[rgb]{0.31,0.60,0.02}{#1}}
\newcommand{\StringTok}[1]{\textcolor[rgb]{0.31,0.60,0.02}{#1}}
\newcommand{\VariableTok}[1]{\textcolor[rgb]{0.00,0.00,0.00}{#1}}
\newcommand{\VerbatimStringTok}[1]{\textcolor[rgb]{0.31,0.60,0.02}{#1}}
\newcommand{\WarningTok}[1]{\textcolor[rgb]{0.56,0.35,0.01}{\textbf{\textit{#1}}}}
\usepackage{graphicx}
\makeatletter
\newsavebox\pandoc@box
\newcommand*\pandocbounded[1]{% scales image to fit in text height/width
  \sbox\pandoc@box{#1}%
  \Gscale@div\@tempa{\textheight}{\dimexpr\ht\pandoc@box+\dp\pandoc@box\relax}%
  \Gscale@div\@tempb{\linewidth}{\wd\pandoc@box}%
  \ifdim\@tempb\p@<\@tempa\p@\let\@tempa\@tempb\fi% select the smaller of both
  \ifdim\@tempa\p@<\p@\scalebox{\@tempa}{\usebox\pandoc@box}%
  \else\usebox{\pandoc@box}%
  \fi%
}
% Set default figure placement to htbp
\def\fps@figure{htbp}
\makeatother
\setlength{\emergencystretch}{3em} % prevent overfull lines
\providecommand{\tightlist}{%
  \setlength{\itemsep}{0pt}\setlength{\parskip}{0pt}}
\usepackage{bookmark}
\IfFileExists{xurl.sty}{\usepackage{xurl}}{} % add URL line breaks if available
\urlstyle{same}
\hypersetup{
  pdftitle={ChIP-Seq Analysis},
  hidelinks,
  pdfcreator={LaTeX via pandoc}}

\title{ChIP-Seq Analysis}
\author{}
\date{\vspace{-2.5em}2026-04-17}

\begin{document}
\maketitle

\subsection{Introduction}\label{introduction}

ChIP-seq (chromatin immunoprecipitation followed by sequencing)
localises protein-DNA binding events genome-wide --- transcription
factor binding sites, histone modifications, chromatin remodellers. The
resulting alignment density forms peaks at the binding loci; sharp peaks
correspond to point-source binders (TFs), and broad domains correspond
to histone-mark territories (H3K27me3, H3K9me3 spanning kilobases). The
standard pipeline is alignment, peak calling with MACS2, annotation with
ChIPseeker, motif enrichment with HOMER or \texttt{motifmatchr}, and
differential binding with DiffBind.

\subsection{Prerequisites}\label{prerequisites}

A working understanding of read alignment (BWA, Bowtie2), peak calling,
and regulatory genomics including transcription factor binding sites and
histone modifications.

\subsection{Theory}\label{theory}

Two peak-calling modes:

\begin{itemize}
\tightlist
\item
  \textbf{Narrow peaks} (TF ChIP, transcription start sites): MACS2
  defaults are tuned for sharp, point-source binding events.
\item
  \textbf{Broad peaks} (H3K27me3, H3K9me3, H3K36me3): use
  \texttt{-\/-broad} to call diffuse domains spanning kilobases.
\end{itemize}

A matched input control (IgG or total chromatin) is essential for
background normalisation; ChIP enrichment is measured as ChIP-vs-input
fold enrichment per peak region. Standard QC includes:

\begin{itemize}
\tightlist
\item
  \textbf{FRiP} (fraction of reads in peaks): \textgreater{} 1 \%
  indicates successful IP; lower values suggest weak antibody or
  insufficient enrichment.
\item
  \textbf{TSS enrichment}: read coverage at transcription start sites;
  high values confirm chromatin accessibility was preserved.
\item
  \textbf{Cross-correlation}: relative strand cross-correlation should
  peak at the fragment length, indicating successful IP.
\end{itemize}

\subsection{Assumptions}\label{assumptions}

Cross-linking and immunoprecipitation are efficient and specific; the
control input matches ChIP depth; duplicates have been removed
(essential for ChIP-seq peak calling).

\subsection{R Implementation}\label{r-implementation}

\begin{Shaded}
\begin{Highlighting}[]
\FunctionTok{library}\NormalTok{(ChIPseeker)}
\FunctionTok{library}\NormalTok{(TxDb.Hsapiens.UCSC.hg19.knownGene)}

\CommentTok{\# Example: annotate a toy peak set}
\FunctionTok{set.seed}\NormalTok{(}\DecValTok{2026}\NormalTok{)}
\NormalTok{peaks }\OtherTok{\textless{}{-}} \FunctionTok{GRanges}\NormalTok{(}
  \AttributeTok{seqnames =} \FunctionTok{rep}\NormalTok{(}\StringTok{"chr1"}\NormalTok{, }\DecValTok{100}\NormalTok{),}
  \AttributeTok{ranges =} \FunctionTok{IRanges}\NormalTok{(}\AttributeTok{start =} \FunctionTok{sort}\NormalTok{(}\FunctionTok{sample}\NormalTok{(}\FloatTok{1e7}\SpecialCharTok{:}\FloatTok{5e7}\NormalTok{, }\DecValTok{100}\NormalTok{)), }\AttributeTok{width =} \DecValTok{400}\NormalTok{)}
\NormalTok{)}

\NormalTok{anno }\OtherTok{\textless{}{-}} \FunctionTok{annotatePeak}\NormalTok{(peaks,}
                     \AttributeTok{TxDb =}\NormalTok{ TxDb.Hsapiens.UCSC.hg19.knownGene,}
                     \AttributeTok{level =} \StringTok{"gene"}\NormalTok{,}
                     \AttributeTok{verbose =} \ConstantTok{FALSE}\NormalTok{)}
\FunctionTok{as.data.frame}\NormalTok{(anno)[}\DecValTok{1}\SpecialCharTok{:}\DecValTok{5}\NormalTok{, }\FunctionTok{c}\NormalTok{(}\StringTok{"seqnames"}\NormalTok{, }\StringTok{"annotation"}\NormalTok{, }\StringTok{"distanceToTSS"}\NormalTok{)]}

\CommentTok{\# Feature{-}type distribution}
\FunctionTok{plotAnnoBar}\NormalTok{(anno)}

\CommentTok{\# (Motif{-}enrichment example would require sequence extraction with}
\CommentTok{\#  BSgenome and motifmatchr::matchMotifs against a JASPAR PWM set.)}
\end{Highlighting}
\end{Shaded}

\subsection{Output \& Results}\label{output-results}

\texttt{annotatePeak()} assigns each peak to its nearest gene and
categorises as promoter, intron, intergenic, etc. \texttt{plotAnnoBar()}
summarises the feature-type distribution; typical TF ChIP-seq shows
\textasciitilde30 \% promoter, \textasciitilde40 \% intron,
\textasciitilde25 \% distal-intergenic.

\subsection{Interpretation}\label{interpretation}

A reporting sentence: ``MACS2 (with input control, \(q < 0.01\)) called
18\{,\}230 high-confidence STAT1 ChIP peaks in IFN-γ-stimulated cells;
42 \% were promoter-proximal, 35 \% intronic, 18 \% distal-intergenic.
De novo motif analysis (HOMER) recovered the canonical STAT1 GAS motif
(TTCN3GAA) at \(p < 10^{-100}\), confirming antibody specificity. FRiP
was 12 \%, TSS enrichment 16, both well within QC standards.'' Always
report FRiP, TSS enrichment, and motif validation.

\subsection{Practical Tips}\label{practical-tips}

\begin{itemize}
\tightlist
\item
  Always use a matched input control (IgG, sonicated input, or
  no-antibody mock); without it, MACS2 falls back to generic background
  assumptions and inflates false positives.
\item
  Check FRiP per sample; values below 1 \% indicate weak IP, incorrect
  peak-calling mode, or sample failure.
\item
  For histone marks (especially H3K27me3, H3K36me3, H3K9me3), use
  \texttt{-\/-broad}; sharp peak callers miss diffuse domains entirely.
\item
  Validate antibody specificity with de novo motif enrichment (HOMER's
  \texttt{findMotifsGenome.pl} or MEME-ChIP); recovering the expected
  motif is the gold standard for ChIP success.
\item
  For differential binding across conditions, use DiffBind which builds
  consensus peak sets and runs DESeq2 or edgeR on peak-level read
  counts.
\item
  For CUT\&RUN or CUT\&Tag (the modern higher-resolution alternatives),
  the same downstream tools apply; specialised peak callers (SEACR)
  handle the lower-background data more accurately.
\end{itemize}

\subsection{R Packages Used}\label{r-packages-used}

\texttt{ChIPseeker} for peak annotation and visualisation;
\texttt{DiffBind} for differential-binding analysis;
\texttt{motifmatchr} for motif enrichment with PWM databases;
\texttt{chromVAR} for chromatin-accessibility-style motif activity
scores; \texttt{ChIPQC} for quality control on ChIP-seq experiments.

\subsection{For Reviewers}\label{for-reviewers}

\textbf{What to look for in a paper using this method.}

\begin{itemize}
\tightlist
\item
  \textbf{Common misapplications.}
\item
  \textbf{Diagnostics that should be reported but often aren't.}
\item
  \textbf{Red flags in tables and figures.}
\item
  \textbf{What to verify.}
\item
  \textbf{What an adequate Methods paragraph must contain.}
\end{itemize}

\subsection{See also --- labs in this
chapter}\label{see-also-labs-in-this-chapter}

\begin{itemize}
\tightlist
\item
  \href{labs/lab-rna-seq-with-deseq2-edger.Rmd}{RNA-seq with DESeq2 /
  edgeR}
\item
  \href{labs/lab-enrichment-analysis-gsea-ora.Rmd}{Enrichment analysis
  (GSEA / ORA)}
\item
  \href{labs/lab-scrna-seq-with-seurat.Rmd}{scRNA-seq with Seurat}
\end{itemize}

\end{document}
